\documentclass[a4paper]{article}
\usepackage{amsmath}
\usepackage{ctex}
\usepackage{tikz}
\usetikzlibrary{intersections,decorations.text,shadings,3d,positioning,patterns}
\definecolor{c1}{RGB}{62,97,127}
\definecolor{c2}{RGB}{104,182,182}
\definecolor{c3}{RGB}{107,190,190}
\definecolor{c4}{RGB}{100,172,174}
\definecolor{c5}{RGB}{95,162,162}
\definecolor{c6}{RGB}{235,242,252}
\usepackage{hyperref}
\newcommand{\zhongsong}{\CJKfontspec{STZhongsong}}%华文中宋
\newcommand{\xinwei}{\CJKfontspec{STXinwei}}%华文新魏
\usepackage{lmodern}
\begin{document}
\thispagestyle{empty}
\begin{tikzpicture}[remember picture,overlay,font=\sffamily\bfseries]
  \fill[opacity=0.2,c6!50](current page.south east)rectangle (current page.north west);%遮住页码
%===================网格线==============================
\draw[help lines,step=0.8cm,opacity=0.15,color=c1](current page.south east)grid(current page.north west);
%=================右上方弧线图像========================
   \draw[ultra thick,c4,name path=big arc] ([xshift=-2mm]current page.north) arc(150:285:11)
 coordinate[pos=0.225] (x0);
 \begin{scope}
  \clip ([xshift=-2mm]current page.north) arc(150:285:11) --(current page.north east);
  \fill[c4!50,opacity=0.25] ([xshift=4.55cm]x0) circle (4.55);
  \fill[c4!50,opacity=0.25] ([xshift=3.4cm]x0) circle (3.4);
  \fill[c4!50,opacity=0.25] ([xshift=2.25cm]x0) circle (2.25);
  \draw[ultra thick,c4!50] (x0) arc(-90:30:6.5);
  \draw[ultra thick,c4] (x0) arc(90:-30:8.75);
  \draw[ultra thick,c4!50,name path=arc1] (x0) arc(90:-90:4.675);
  \draw[ultra thick,c4!50] (x0) arc(90:-90:2.875);
  \path[name intersections={of=big arc and arc1,by=x1}];
  \draw[ultra thick,c4,name path=arc2] (x1) arc(135:-20:4.75);
  \draw[ultra thick,c4!50] (x1) arc(135:-20:8.75);
  \path[name intersections={of=big arc and arc2,by={aux,x2}}];
  \draw[ultra thick,c4!50] (x2) arc(180:50:2.25);
 \end{scope}
 %================日期==================================
 \path[decoration={text along path,text color=c4,
                 raise = -2.8ex,
                 text = {|\sffamily\bfseries|\today},
                 text align = center
             },
             decorate
         ] ([xshift=-2mm]current page.north) arc(150:245:11);

%===================坐标轴==============================
\node[opacity=0.5,color=c1] at([xshift=4cm,yshift=-4cm]current page.north
  west){\tikz{\draw[-stealth](-0.3,0)--(0,0)node[below left]{$O$}--(3.6,0)node[below=2pt]{$x$};
\draw[-stealth,color=c1](0,-0.42)--(0,2)node[left]{$y$};
\draw(0,1.3)node[left]{$A$}..controls(1.4,0.6)and(2.2,2.4)..(2.8,1.7)node[right]{$B$}
--(2.8,0)node[below]{$C$};\node at(0.65,1.4){$f(x)$};}};
%=================标题信息=========================
\node[text=c1,anchor=west,scale=4,inner sep=0pt,font=\zhongsong] at
 ([xshift=-10.0cm,yshift=-1.0cm]current page.center) {高中数学题库};

\node[text=c1,anchor=west,scale=2.5,inner sep=0pt,font=\zhongsong] at
 ([xshift=-9.85cm,yshift=-2.65cm]current page.center) {\fontencoding{T1}\fontfamily{ptm}\bfseries\selectfont High School Mathematics Problems};

\node[text=c1,anchor=west,scale=2.0,inner sep=0pt,font=\zhongsong] at
 ([xshift=-10.15cm,yshift=-3.75cm]current page.center) {（Version 0.1）};

\draw[ultra thick,c1!50]([xshift=0.3cm,yshift=-5cm]current page.west)--([xshift=-0.3cm,yshift=-5cm]current page.east);

\node[text=c1,anchor=west,scale=2,inner sep=0pt,font=\heiti] at
 ([xshift=-9.75cm,yshift=-6cm]current page.center) {槿灵兮};

\node[opacity=0.5,scale=1.8,color=c1!80] at ([xshift=-3.15cm,yshift=-8.0cm]current page.center) {$S_{\triangle BPC}\overrightarrow{PA} + S_{\triangle CPA}\overrightarrow{PB} + S_{\triangle APB}\overrightarrow{PC} = \overrightarrow{0}$};

\node[opacity=0.5,scale=1.8,color=c1!80] at ([xshift=4.30cm,yshift=-10.0cm]current page.center) {$x_1 y_1 - x_2 y_2 \geq \sqrt{x_1^2 - x_2^2} \sqrt{y_1^2 - y_2^2}$};

%==============出版社信息=============================
\node[text=c1,scale=2,inner sep=0pt,font=\heiti] at ([yshift=2cm]current page.south){数圈MathCyclus};

\end{tikzpicture}

\end{document}
% 创作不易，转载请注明出处，谢谢合作！